% ===========================================================================
% Source material: docs/architecture.md sections 1 and 2
% The template grades "actors/stakeholders' requirement gathering with
% justification", so every requirement below needs a stated reason.
% ===========================================================================

\section{Requirement Gathering}

\subsection{Actors and stakeholders}

\begin{table}[htbp]
  \centering
  \caption{Actors, what the system permits them, and what it structurally prevents.}
  \label{tab:actors}
  \small
  \begin{tabular}{@{}L{2.8cm}L{6.0cm}L{6.4cm}@{}}
    \toprule
    \textbf{Actor} & \textbf{Can do}                                                                  & \textbf{Cannot do} \\
    \midrule
    Election Authority
                   & Create an election, set candidates and windows, publish metadata, cancel before voting
                   & Cast or alter a ballot; change the root once registration closes; stop the count \\
    Registrar
                   & Check eligibility off-chain, collect commitments, publish the roll's root, size and leaves
                   & Learn from chain data how anyone voted; change a leaf after the root freezes    \\
    Voter
                   & Register a commitment, cast one ballot, verify that it landed
                   & Vote twice; vote from outside the committed set                                 \\
    Candidate
                   & Be listed on a ballot, publish a manifesto
                   & Anything privileged                                                             \\
    Auditor / Observer
                   & Re-derive the tally, verify every proof, recompute the root from the leaves
                   & Alter anything                                                                  \\
    Relayer
                   & Pay gas for a voter; read the choice in every ballot; refuse to forward one
                   & Forge, modify or replay a ballot; learn a voter's identity from the payload   \\
    \bottomrule
\end{tabular}
\end{table}

\paragraph{Election Authority.}
It cannot finalize: closing voting and fixing the result are permissionless
once their timestamps pass, so it cannot stall a poll it is losing. Every act
reserved to it precedes the first ballot and emits an event. More power would
let it enrol itself; with less, no election could be convened (D5).

\paragraph{Registrar.}
Eligibility is an off-chain fact whose evidence must never reach a public
ledger, so the registrar checks it (institute email plus a one-time password)
and publishes only the root, roll size and leaf set identifier. The booth
derives its own path from that set, so the registrar is absent from the ballot
path (FR15). The root freezes when registration closes, since a registrar able
to append leaves mid-vote could manufacture ballots.

\paragraph{Voter.}
The identity secret is generated in the browser; only its commitment reaches
the registrar, except under Section~6's proving fallback, which merges
registration, proving and relaying into one trust domain. A voter who cannot
check that their ballot landed has only moved trust to a contract; a provable
ballot is also a coercible one. A lost secret cannot be replaced once the root
freezes, so the booth exports a copy and re-registration closes with the
window.

\paragraph{Candidate.}
The candidate holds no privileged capability: every privileged operation
affects the outcome they compete for. Nothing they need requires one, since the
roster and manifesto identifier are on-chain, so substitution is detectable,
and the tally is public.

\paragraph{Auditor and observer.}
Auditing needs no registration or approval: FR12 makes the ballots, tally, root
and roll size public, and NFR11 covers the roll check. An audit through this
project's backend would verify the operator, not the chain; D5 recomputes the
tally with it off.

\paragraph{Relayer.}
The relayer exists so a voter need not acquire cryptocurrency (FR10). It
cannot alter a ballot: the nullifier and candidate signal are bound inside the
proof, and the relayed request is only accounting. The booth builds its own
path, so the relayer is the only service this project must keep answering
during voting. That it reads the choice while also deciding whether to forward
is conceded below as one attack.

\subsection{Functional requirements}

\begin{longtable}{@{}>{\RaggedRight\arraybackslash\small}p{1.3cm}>{\RaggedRight\arraybackslash\small}p{8.1cm}>{\RaggedRight\arraybackslash\small}p{7.9cm}@{}}
  \caption{Functional requirements, each derived from an actor in Table~\ref{tab:actors}.}\\
  \toprule
  \textbf{ID} & \textbf{The system shall} & \textbf{Why, and for which actor} \\
  \midrule
  \endhead
  \bottomrule
  \endlastfoot
  FR1 & Let the Authority create an election through \texttt{ElectionFactory} with the roster, two windows and metadata identifier. & Only the Authority may define what is decided; on-chain first, so changes stay visible. \\
  \addlinespace[2pt]
  FR2 & Permit candidate changes only in \texttt{Created} or \texttt{Registration}. & Errors must be fixable, but a roster changing mid-vote alters the meaning of ballots already cast. \\
  \addlinespace[2pt]
  FR3 & Publish the ballot description and each manifesto to IPFS, storing the content identifier on-chain. & Voters need a readable statement; the committed identifier makes substitution detectable. \\
  \addlinespace[2pt]
  FR4 & Let the Registrar check eligibility off-chain and accept each identity commitment, writing no personal data on-chain. & Only the Registrar can establish this off-chain fact, whose evidence must never reach a public ledger. \\
  \addlinespace[2pt]
  FR5 & Have the Registrar write the root, roll size and leaf set identifier to \texttt{VoterRegistry} in one transaction when registration closes, publish the leaf set to IPFS, and reject later root changes. & Auditors cannot recompute an unpublished root, the booth cannot build a path without the leaves, and freezing blocks mid-vote appends. \\
  \addlinespace[2pt]
  FR6 & Require an explicit, irreversible transaction per phase transition, rejected before its block timestamp; reserve configuration and cancellation to the Authority before voting; any account may close and finalize. & Stops the Authority reopening a poll it dislikes or cancelling after ballots exist, and stops it stalling by never calling. \\
  \addlinespace[2pt]
  FR7 & Accept a ballot only with a zero-knowledge membership proof against the published root, a nullifier and a signal encoding the choice, verified before the tally moves. & The Voter must vote without revealing who they are; others need each ballot provably from the roll. \\
  \addlinespace[2pt]
  FR8 & Record every accepted nullifier, revert on a repeat, and revert once spent nullifiers would exceed the roll size stored beside the root. & One ballot per voter without recording addresses, which would defeat secrecy; the cap makes capacity a contract invariant. \\
  \addlinespace[2pt]
  FR9 & Let a Voter holding gas call \texttt{castVote} directly, without the Relayer. & A Relayer that reads the choice and refuses to forward could otherwise alter the result unobserved. \\
  \addlinespace[2pt]
  FR10 & Let the Relayer accept a signed ballot, verify its proof off-chain, pre-check its nullifier, return a signed receipt, and forward it through \texttt{VoteForwarder} under ERC-2771, paying the gas. & The Voter must hold no cryptocurrency (E2); pre-verification stops junk proofs draining the Relayer, and the receipt makes a dropped ballot evidence. \\
  \addlinespace[2pt]
  FR11 & Increment tallies through \texttt{TallyLib} in the transaction accepting a ballot, and let any account finalize once voting closes; the result is then immutable. & No actor may adjust the count between the last ballot and the announcement, or leave it unannounced. \\
  \addlinespace[2pt]
  FR12 & Make every ballot event, the tally, the root, the roll size and the result readable from chain data by anyone, needing nothing this project runs. & The Auditor is permissionless; an audit needing the operator's cooperation audits the operator. \\
  \addlinespace[2pt]
  FR13 & Return the Voter their ballot's transaction identifier, and warn them and offer direct submission when none arrives before a stated deadline. & The Voter's only direct evidence of acceptance; a ballot never forwarded should leave a signed promise, not silence. \\
  \addlinespace[2pt]
  FR14 & Have the Registrar return a receipt signed over the accepted commitment and the election identifier at registration. & The Voter otherwise holds no evidence of registering; with it, exclusion from the leaf set can be contradicted. \\
  \addlinespace[2pt]
  FR15 & Have the booth download the published leaf set, confirm the Voter's commitment is present, rebuild the tree and derive the authentication path locally, asking the Registrar for nothing at vote time. & A path from the Registrar would reveal which leaf was voting, and the Relayer sees the choice, so the two records deanonymize the Voter. It also drops the Registrar as a liveness dependency, and the presence check catches substitution. \\
\end{longtable}

\subsection{Non-functional requirements}

\begin{longtable}{@{}>{\RaggedRight\arraybackslash\small}p{1.5cm}>{\RaggedRight\arraybackslash\small}p{10.4cm}>{\RaggedRight\arraybackslash\small}p{5.4cm}@{}}
  \caption{Non-functional requirements and the Section~5 parameter that will test each one.}\\
  \toprule
  \textbf{ID} & \textbf{Requirement} & \textbf{Tested by} \\
  \midrule
  \endhead
  \bottomrule
  \endlastfoot
  NFR1 & \emph{Ballot secrecy, conditional.} No chain data ties a ballot to a commitment or registration record; the sender is an ephemeral key never checked against the roll. Two qualifications belong to the requirement: it is anonymity within the roll, not secrecy of ballot contents; and it is unconditional only against the chain, holding against the operator while the Relayer runs separately from the Registrar, or the voter submits directly. & D3, a deanonymization attempt on full chain data, reporting the anonymity set left \\
  \addlinespace[2pt]
  NFR2 & \emph{Eligibility enforcement.} A ballot is accepted only on proof of membership of the published root, subject to the trusted setup assumption below. & D2: a proof against the wrong root must revert \\
  \addlinespace[2pt]
  NFR3 & \emph{One ballot per voter.} A second ballot from the same identity secret is rejected, and accepted ballots never exceed the published roll size. & D1, with fuzzing over nullifier collisions \\
  \addlinespace[2pt]
  NFR4 & \emph{Individual verifiability.} A voter can locate their own ballot in chain data using only what their client holds, and confirm their commitment in the leaves beforehand. & D4 \\
  \addlinespace[2pt]
  NFR5 & \emph{Universal verifiability.} Any party can recompute the tally from raw chain logs and reach the published result. & D4, with this project's backend off \\
  \addlinespace[2pt]
  NFR6 & \emph{Integrity of the count.} Tallies move only on the ballot-acceptance path, and no role, the Authority included, can alter a tally, stop the count or withhold finalization. & D5 \\
  \addlinespace[2pt]
  NFR7 & \emph{Cost of running an election.} A ballot including proof verification stays under 400,000 gas and under INR 0.50 on Polygon Amoy; a new election deploys as a clone. & A1 and A2 \\
  \addlinespace[2pt]
  NFR8 & \emph{End-to-end latency.} The 95th percentile from click to on-chain inclusion on Polygon Amoy stays under 10 seconds, of which browser proving is under 5 seconds on a mid-range phone. & B1 and B2 \\
  \addlinespace[2pt]
  NFR9 & \emph{Scalability in the electorate.} Gas per operation is constant in the number of enrolled voters once a ballot carries FR7's proof, since neither the root nor the proof grows with the roll. The plain Merkle-proof baseline differs at the ballot, whose authentication path grows with tree depth. Browser proving time grows with the logarithm of the roll. The claim is per operation: election cost stays linear in ballots cast. & C1 and C2, with C3 for the concurrency ceiling \\
  \addlinespace[2pt]
  NFR10 & \emph{No cryptocurrency held by the voter.} Registration and voting complete while the voter holds a zero balance. & E2, with E1 for actions \\
  \addlinespace[2pt]
  NFR11 & \emph{Auditability of the voter roll.} The root, roll size and leaf set identifier are on-chain and the leaves published, so anyone can recompute the root and compare it, and a voter can find their commitment. Two limits belong to the requirement: the check costs work linear in the roll, paid by each auditor separately; and it shows only that the leaves produce the root, never that they are real people. & D4 extended to the root \\
  \addlinespace[2pt]
  NFR12 & \emph{Availability of the public result path.} The result stays readable and verifiable while every service this project runs is offline. & D4 \\
  \addlinespace[2pt]
  NFR13 & \emph{Operability by a non-specialist.} A voter with no cryptocurrency experience casts a ballot unaided, read as more than eight of ten E3 participants finishing without help. & E3 for completion, E1 capping actions at four \\
\end{longtable}

NFR2, NFR3 and NFR9 are load-bearing: ineligible or duplicate ballots make the
number meaningless, and NFR9 is the decision the architecture rests on, with C2
the half that can genuinely fail. Every target is a hypothesis, to be reported
as measured in Evaluation 2, including where missed.

\subsection{Trust model and explicit non-goals}

Groth16 needs a circuit-specific trusted setup, so eligibility rests on the
phase-two ceremony for the Semaphore circuit; this project uses the published
artifacts and verifies their hashes against the Semaphore repository rather
than running one. If every contributor colluded, forged membership proofs would
verify and nothing in the proofs would reveal it. FR8's on-chain cap against
the declared roll size does not depend on the ceremony.

Six guarantees, NFR1 to NFR6 as claims, each paired in Section~5 with a test.

\begin{itemize}
  \item \emph{Eligibility.} Only an identity in the published root can produce
        a proof the contract accepts, given the setup assumption.
  \item \emph{One voter, one ballot.} Nullifier uniqueness, without tracking
        addresses, plus the contract's roll-size cap.
  \item \emph{Ballot secrecy.} The proof reveals nothing beyond membership. It
        is an anonymity set, worth nothing for a lone ballot in a quiet hour,
        and unconditional only against the chain: against the operator it needs
        the relayer run separately from the registrar, or direct submission.
  \item \emph{Individual verifiability.} A voter can find their own ballot
        transaction, and their commitment in the leaves.
  \item \emph{Universal verifiability.} Anyone can recompute the tally from
        public chain data, unaided.
  \item \emph{Integrity of the count.} The tally moves only on the
        ballot-acceptance path, and closing and finalizing are permissionless,
        so no operator can adjust or withhold it.
\end{itemize}

Six properties \projectname{} does not guarantee:

\begin{itemize}
  \item \emph{Coercion resistance, and the sale of a credential.} The identity
        secret is a bearer credential, so handing it over transfers the vote
        irrevocably: the nullifier is deterministic and there is no key
        rotation. The buyer verifies delivery on-chain, remotely and in bulk,
        so here the design is worse than a paper ballot, which lets a coerced
        voter lie. The remedy, MACI-style key rotation, needs a coordinator and
        circuit work this schedule cannot fit.
  \item \emph{Sybil resistance at the registrar.} A dishonest registrar can
        enrol identities matching no real person, and publishing the roll size
        does not make that detectable: registration is voluntary, so nobody
        knows the true figure. Publication buys a bound on capacity, not
        detection: the number is fixed and public before any ballot exists, and
        \texttt{castVote} refuses ballots past it. A registrar can still enrol
        phantom leaves early and spend them late for whichever candidate the
        running tally needs. The remedy is a split registrar, or observers
        co-signing the leaf set; neither is built.
  \item \emph{Roll correctness.} Nothing forces the frozen roll to match who
        actually registered: a registrar can drop or substitute a commitment
        and the leaf set still reproduces the published root, since recomputing
        a root proves only that the leaves match it. FR14 and FR15 make that
        provable rather than invisible, and the window to act closes when
        voting opens. Neither makes the registrar honest.
  \item \emph{Privacy from the relayer, and selective censorship.} The relayer
        reads the candidate choice in every ballot and sees the source IP. The
        payload names nobody, but one operator running the registrar and the
        relayer can correlate its own two logs, so secrecy from the operator is
        conditional on separation. With the freedom to refuse, this is one
        attack: dropping only the ballots naming one candidate alters the
        result, and turnout below the roll size hides it. FR10's receipt and
        FR13's deadline warning bound it, and FR9 only helps a voter holding
        gas; removing it needs encrypted ballots or several independent
        relayers.
  \item \emph{Network-level anonymity.} A relayer or RPC provider sees the IP a
        ballot arrives from, and an operator that also ran registration holds
        the address that person registered from. The relayer strips what
        metadata it can and keeps no request logs, but that is a promise by the
        distrusted party, not a property. A real fix means Tor or a mixnet, out
        of scope here.
  \item \emph{Secrecy of the running tally.} Ballots are counted on-chain as
        they arrive, so a partial result is visible while voting is open and
        can sway voters yet to cast. It also erodes ballot secrecy: timestamped
        \texttt{VoteCast} events with a public counter let anyone who knows
        when a person voted narrow the anonymity set. Hiding it needs
        homomorphic or commit-reveal tallying, which adds a phase the Section~4
        lifecycle does not model.
\end{itemize}

Each follows from a scope decision; the remedies above are future work.
